% === INSERIRE DOPO l'ultima \section{} già presente nel documento ===

% ============================================================
\section{Hashing e Tabelle Hash}
% ============================================================

\begin{softnote}
	Le tabelle hash sono strutture dati efficienti per implementare dizionari
	che supportano le sole operazioni \textsc{Insert}, \textsc{Search} e
	\textsc{Delete}. Analogamente a un array indicizzato da interi, una tabella
	hash indicizza gli elementi tramite la \emph{chiave} dell'elemento stesso.
	Nel caso medio, sotto ipotesi ragionevoli, il tempo di ricerca è $O(1)$,
	anche se nel caso peggiore è $O(n)$.
\end{softnote}

% ------------------------------------------------------------
\subsection{Tavole a indirizzamento diretto}
% ------------------------------------------------------------

\begin{definizione}[Tavola a indirizzamento diretto]
	Sia $U$ l'insieme di tutti i possibili valori delle chiavi. Se $|U|$ non è
	troppo grande e nessun elemento ha chiave uguale a un altro, è possibile
	convertire ogni chiave in un intero e usare un array $T$ di lunghezza $|U|$
	per memorizzare i riferimenti agli elementi. La funzione di indicizzazione è:
	\[
	I(k) = \sum_{i=0}^{n-1} \bigl(k[i] - \texttt{'A'}\bigr) \cdot 26^i
	\]
	(nell'esempio con chiavi alfanumeriche di 2 caratteri in
	$[\texttt{A},\ldots,\texttt{Z}]$, $|U| = 26^2 = 676$).
\end{definizione}

\begin{hardnote}
	L'indirizzamento diretto è praticabile solo quando $|U|$ è piccolo. Se
	$n \ll |U|$ (poche chiavi usate su un universo enorme) si spreca molta
	memoria: si deve allocare un array di dimensione $|U|$ anche se si
	memorizzano solo $n$ elementi.
\end{hardnote}

% ------------------------------------------------------------
\subsection{Tabelle hash}
% ------------------------------------------------------------

Quando $|U|$ è troppo grande, si usa una tabella di dimensione $m \approx n$
(con $m \ll |U|$) e una \textbf{funzione hash}
\[
h : U \to \{0, \ldots, m-1\}
\]
che mappa la chiave $k$ nell'indice $h(k)$ (\emph{valore hash} di $k$).
Poiché $m \ll |U|$, $h$ non può essere iniettiva: possono quindi verificarsi
\textbf{collisioni}.

\begin{definizione}[Funzione hash perfetta]
	Una funzione hash $h$ è \emph{perfetta} per un insieme di chiavi $S$ se
	non produce collisioni sulle sole chiavi di $S$, ovvero se $h$ è iniettiva
	su $S$.
\end{definizione}

% ------------------------------------------------------------
\subsection{Funzioni hash: scelta e qualità}
% ------------------------------------------------------------

\begin{softnote}
	Una funzione hash ideale determina ogni valore $h(k)$ in modo casuale e
	indipendente nell'insieme $\{0,\ldots,m-1\}$, restituendo sempre lo stesso
	valore per la stessa chiave (\emph{funzione di hash uniforme e indipendente}
	o \emph{oracolo casuale}). Tale funzione è solo teorica; in pratica si usano
	funzioni che la approssimano.
\end{softnote}

\subsubsection{Conversione verso gli interi}

Quando le chiavi non sono interi, si converte prima la chiave in un intero:

\begin{itemize}
	\item \textbf{Caratteri Unicode (Java)}: base $65536$ (16 bit per
	\texttt{char}):
	\[
	k' = \sum_{i} c_i \cdot 65536^i
	\]
	\item \textbf{Caratteri Latin-1}: base $256$.
	\item \textbf{Subset ristretto}: base ulteriormente ridotta riassegnando
	gli ordinali a partire da $0$.
\end{itemize}

\begin{nota}
	Per stringhe di anche soli 2 caratteri, la somma $k'$ supera il limite di
	un \texttt{int} Java (32 bit). Il risultato viene troncato a $w = 32$ bit,
	equivalente a calcolare $k' \bmod 2^{32}$. La conversione perde biettività
	(possibili collisioni), ma questo è accettabile: anche $h$ stessa non è
	iniettiva. L'obiettivo è semplicemente \emph{mescolare} i caratteri in un
	intero prima di applicare $h$.
\end{nota}

% ------------------------------------------------------------
\subsection{Hashing statico}
% ------------------------------------------------------------

\subsubsection{Metodo della divisione}

\begin{definizione}[Metodo della divisione]
	\[
	h(k) = k \bmod m
	\]
	Richiede una sola operazione. La scelta di $m$ è critica: si preferisce un
	\textbf{numero primo non vicino a potenze di 2}, oppure un numero con
	fattori primi $> 20$ (Lum, 1971).
\end{definizione}

\begin{hardnote}
	Se $m = 2^p$, allora $h(k)$ dipende solo dagli $p$ bit meno significativi
	di $k$, ignorando completamente la parte alta della chiave. Ad esempio, con
	$m = 2^{12} = 4096$:
	\[
	h(40960) = h(45056) = h(49152) = h(53248) = 0
	\]
	poiché questi valori hanno tutti i 12 bit meno significativi pari a zero.
\end{hardnote}

\subsubsection{Metodo della moltiplicazione}

\begin{definizione}[Metodo della moltiplicazione]
	Sia $0 < A < 1$ una costante. Il valore hash è:
	\[
	h(k) = \bigl\lfloor m \cdot \bigl((k \cdot A) \bmod 1\bigr) \bigr\rfloor
	\]
	dove $(k \cdot A) \bmod 1$ è la parte frazionaria di $k \cdot A$. Il valore
	$m$ non è critico; si usa tipicamente $m = 2^p$.
\end{definizione}

\begin{nota}
	Donald Knuth suggerisce $A \approx \dfrac{\sqrt{5}-1}{2} \approx 0.6180$
	(sezione aurea).
\end{nota}

\subsubsection{Metodo moltiplica-e-scorri}

Il metodo della moltiplicazione può essere reso più efficiente fissando
$m = 2^l$ con $0 < l \le w$ ($w$ = numero di bit della parola macchina).
Si pone $a = A \cdot 2^w$ (intero dispari), e si calcola:
\[
h_a(k) = \bigl((k \cdot a) \bmod 2^w\bigr) \mathbin{>>>} (w - l)
\]
dove \texttt{>>>} è lo scorrimento logico a destra (inserisce zeri a sinistra,
a differenza di \texttt{>>} che mantiene il bit di segno). Si estraggono così
gli $l$ bit più significativi di $r_0 = (k \cdot a) \bmod 2^w$.

\begin{esempio}[Moltiplica-e-scorri con $w=8$, $m=16$]
	Dati $w = 8$, $m = 2^4 = 16$ (dunque $l = 4$), $k = 202$
	($\mathtt{11001010}_2$), $a = 173$ ($\mathtt{10101101}_2$):
	\[
	k \cdot a = 202 \cdot 173 = 34946, \quad
	r_0 = 34946 \bmod 256 = 130 = \mathtt{10000010}_2
	\]
	\[
	h_a(k) = \mathtt{10000010}_2 \mathbin{>>>} 4 = \mathtt{00001000}_2 = 8
	\]
	La chiave viene quindi assegnata alla cella $8 \in \{0,\ldots,15\}$.
\end{esempio}

% ------------------------------------------------------------
\subsection{Hashing randomizzato}
% ------------------------------------------------------------

\begin{softnote}
	Con una funzione hash statica, un avversario può costruire $n$ chiavi che
	collidono tutte nella stessa cella, portando a $n-1$ collisioni (caso
	peggiore sistematico). L'hashing randomizzato risolve questo problema
	scegliendo casualmente la funzione hash \emph{all'inizio di ogni esecuzione}
	da una famiglia opportuna: nessun input può essere sistematicamente il caso
	peggiore.
\end{softnote}

\begin{definizione}[Hashing universale]
	Sia $\mathcal{H}$ una famiglia finita di funzioni hash con dominio $U$ e
	codominio $\{0,1,\ldots,m-1\}$. $\mathcal{H}$ è \emph{universale} se per
	ogni coppia di chiavi distinte $k_i, k_j \in U$:
	\[
	\Pr_{h \in \mathcal{H}}\bigl\{h(k_i) = h(k_j)\bigr\} \le \frac{1}{m}
	\]
	Equivalentemente, il numero di funzioni $h \in \mathcal{H}$ che fanno
	collidere $k_i$ e $k_j$ è al più $|\mathcal{H}|/m$.
\end{definizione}

\begin{definizione}[Classificazione delle famiglie]
	Sia $\mathcal{H}$ una famiglia finita di funzioni hash
	($U \to \{0,\ldots,m-1\}$) e $h \in \mathcal{H}$ scelta casualmente:
	\begin{itemize}
		\item $\mathcal{H}$ è \textbf{uniforme} se $\forall k \in U$ e
		$\forall q \in \{0,\ldots,m-1\}$: $\Pr\{h(k)=q\} = 1/m$.
		\item $\mathcal{H}$ è $\varepsilon$\textbf{-universale} se
		$\forall k_i \ne k_j$:
		$\Pr\{h(k_i)=h(k_j)\} \le \varepsilon$.
		\item $\mathcal{H}$ è \textbf{universale} se è
		$\varepsilon$-universale con $\varepsilon = 1/m$.
	\end{itemize}
	Una famiglia uniforme e indipendente è anche universale (ma non viceversa).
\end{definizione}

\paragraph{Famiglia basata sulla teoria dei numeri.}
Sia $p$ un numero primo tale che ogni chiave $k \in U$ sia associabile a un
intero in $[0,\ldots,p-1]$, con $p > m$. Siano
$\mathbb{Z}_p = \{0,\ldots,p-1\}$ e $\mathbb{Z}_p^* = \{1,\ldots,p-1\}$.
Per ogni $a \in \mathbb{Z}_p^*$ e $b \in \mathbb{Z}_p$ si definisce:
\[
h_{ab}(k) = \bigl((ak + b) \bmod p\bigr) \bmod m
\]
La famiglia $\mathcal{H}_{pm} = \{h_{ab} \mid a \in \mathbb{Z}_p^*,\,
b \in \mathbb{Z}_p\}$ è universale, con $|\mathcal{H}_{pm}| = p(p-1)$ e
$m$ arbitrario.

\paragraph{Famiglia $2/m$-universale basata su moltiplica-e-scorri.}
La famiglia:
\[
\mathcal{H} = \bigl\{h_a(\cdot) \;\big|\; 1 \le a < 2^w,\;
a \text{ dispari}\bigr\}
\]
è $2/m$-universale (con $m = 2^l$). La velocità di calcolo di $h_a(k)$
(una moltiplicazione e uno scorrimento) compensa la probabilità di collisione
$2/m$ invece di $1/m$.

\begin{hardnote}
	Nella versione italiana di [CLRS] sono presenti due errori nella definizione
	di questa famiglia: verificare sempre la definizione originale (edizione
	inglese).
\end{hardnote}

% ------------------------------------------------------------
\subsection{Gestione delle collisioni}
% ------------------------------------------------------------

\subsubsection{Risoluzione mediante concatenamento}

\begin{softnote}
	Nel concatenamento (\emph{chaining}), ogni cella $T[j]$ della tabella è la
	testa di una lista concatenata. Tutti gli elementi con lo stesso valore hash
	finiscono nella stessa lista.
\end{softnote}

Le operazioni sono:
\begin{itemize}
	\item \textsc{Insert}$(o)$: inserisce $o$ in testa alla lista
	$T[h(o.\mathit{key})]$ --- $\Theta(1)$.
	\item \textsc{Search}$(k)$: scorre la lista $T[h(k)]$ cercando la chiave
	$k$.
	\item \textsc{Delete}$(o)$: rimuove $o$ dalla lista
	$T[h(o.\mathit{key})]$; poiché $o$ contiene i puntatori
	\texttt{prev}/\texttt{next}, l'operazione è $\Theta(1)$.
\end{itemize}

\begin{definizione}[Fattore di carico]
	Il \emph{fattore di carico} (load factor) è:
	\[
	\alpha = \frac{n}{m}
	\]
	dove $n$ è il numero di elementi memorizzati e $m$ la dimensione della
	tabella. Nel concatenamento $\alpha > 1$ è accettabile (preferibilmente
	$\alpha \approx 1$).
\end{definizione}

\begin{teorema}[Costo medio --- ricerca senza successo]
	In una tabella hash con concatenamento, sotto l'ipotesi di hashing uniforme
	e indipendente, una ricerca senza successo richiede in media tempo
	$\Theta(1 + \alpha)$.
\end{teorema}

\begin{proof}
	Ogni lista delle $m$ celle ha lunghezza media $\alpha$. La ricerca di una
	chiave non presente richiede $\Theta(1)$ per calcolare il valore hash e
	$\Theta(\alpha)$ per scorrere la lista.
\end{proof}

\begin{teorema}[Costo medio --- ricerca con successo]
	Sotto le stesse ipotesi, una ricerca con successo richiede in media tempo
	$\Theta(1 + \alpha)$.
\end{teorema}

\begin{proof}
	Sia $x_i$ l'$i$-esimo elemento inserito. Per ogni coppia $(i,j)$ con $j > i$
	e ogni cella $q$, sia $X_{ijq}$ la variabile indicatrice dell'evento
	\{\,si cerca $x_i$, $h(k_i)=q$ e $h(k_j)=q$\,\}.
	Sotto hashing uniforme e indipendente:
	\[
	\Pr\{X_{ijq}=1\} = \frac{1}{n} \cdot \frac{1}{m^2}
	\quad\Rightarrow\quad E[X_{ijq}] = \frac{1}{nm^2}
	\]
	Il numero atteso di elementi che precedono l'elemento cercato nella lista
	vale:
	\[
	E[1 + Z]
	= 1 + \sum_{j=1}^{n}\sum_{q=0}^{m-1}\sum_{i=1}^{j-1} \frac{1}{nm^2}
	= 1 + \frac{\alpha}{2} - \frac{\alpha}{2n}
	\]
	Aggiungendo il costo $\Theta(1)$ per il calcolo dell'hash, il costo totale
	è $\Theta\!\left(2 + \dfrac{\alpha}{2} - \dfrac{\alpha}{2n}\right)
	= \Theta(1 + \alpha)$.
\end{proof}

\begin{table}[ht]
	\centering
	\caption{Prestazioni della tabella hash con concatenamento}
	\begin{tabular}{lcc}
		\toprule
		\textbf{Operazione} & \textbf{Caso peggiore} & \textbf{Caso medio}  \\
		\midrule
		\textsc{Insert}     & $\Theta(1)$             & $\Theta(1)$           \\
		\textsc{Search}     & $\Theta(n)$             & $\Theta(1 + \alpha)$  \\
		\textsc{Delete}     & $\Theta(1)$             & $\Theta(1)$           \\
		\bottomrule
	\end{tabular}
\end{table}

\begin{nota}
	Se $n = O(m)$ (il numero di elementi è proporzionale al numero di celle),
	allora $\alpha = n/m = O(1)$ e nel caso medio tutte le operazioni costano
	$O(1)$.
\end{nota}

% ------------------------------------------------------------
\subsubsection{Risoluzione mediante indirizzamento aperto}
% ------------------------------------------------------------

\begin{softnote}
	Nell'indirizzamento aperto (\emph{open addressing}), tutti gli elementi sono
	memorizzati nella tabella stessa (niente liste esterne). In caso di
	collisione si cerca un'altra cella libera tramite una sequenza di
	\emph{ispezioni} (\emph{probing}). La funzione hash viene estesa come:
	\[
	h : U \times \{0,1,\ldots,m-1\} \to \{0,1,\ldots,m-1\}
	\]
	La funzione hash non estesa è denotata $h_0(\cdot)$. Poiché tutti gli slot
	sono nella tabella, il fattore di carico deve soddisfare $\alpha < 1$.
\end{softnote}

\begin{algorithm}[H]
	\caption{\textsc{Insert}$(T, o)$}
	\KwIn{Tabella hash $T$; oggetto $o$ da inserire (non già presente).}
	\KwOut{Oggetto inserito, oppure errore ``Overflow''.}
	$i \leftarrow 0$\;
	\While{$i \neq m$}{
		$j \leftarrow h(o.\mathit{key},\, i)$\;
		\If{$T[j] = \mathtt{NULL}$ \KwOr $T[j] = \mathtt{DELETE}$}{
			$T[j] \leftarrow o$\;
			\KwRet{``Inserito''}\;
		}
		$i \leftarrow i + 1$\;
	}
	\KwRet{``Overflow''}\;
\end{algorithm}

\begin{algorithm}[H]
	\caption{\textsc{Search}$(T, k)$}
	\KwIn{Tabella hash $T$; chiave $k$ da cercare.}
	\KwOut{Indice di $k$ in $T$ se presente, \texttt{NULL} altrimenti.}
	$i \leftarrow 0$\;
	\While{$i \neq m$}{
		$j \leftarrow h(k,\, i)$\;
		\lIf{$T[j] = \mathtt{NULL}$}{\KwRet{\texttt{NULL}}}
		\lIf{$T[j] \neq \mathtt{DELETE}$ \KwAnd $T[j].\mathit{key} = k$}{\KwRet{$j$}}
		$i \leftarrow i + 1$\;
	}
	\KwRet{\texttt{NULL}}\;
\end{algorithm}

\begin{algorithm}[H]
	\caption{\textsc{Delete}$(T, i)$}
	\KwIn{Tabella hash $T$; indice $i$ della cella da cancellare.}
	$T[i] \leftarrow \mathtt{DELETE}$ \tcp*{marcatore tombstone}
\end{algorithm}

\begin{hardnote}
	Il marcatore \texttt{DELETE} (\emph{tombstone}) è necessario per non
	interrompere le sequenze di ispezione. Se si sostituisse con \texttt{NULL},
	la \textsc{Search} si fermerebbe su quella cella producendo falsi negativi.
	\textsc{Search} non si ferma su \texttt{DELETE} e continua l'ispezione;
	\textsc{Insert} può riutilizzare una cella \texttt{DELETE}. Con molte
	cancellazioni le ispezioni peggiorano: si risolve con un \textbf{rehash}
	periodico (ricostruzione della tabella con i soli elementi presenti).
\end{hardnote}

\paragraph{Analisi (ipotesi).}
L'analisi classica [CLRS] assume \emph{hashing con permutazioni indipendenti e
	uniformi}: per ogni chiave $k$, la sequenza
$\langle h(k,0), \ldots, h(k,m-1)\rangle$ è una permutazione casuale uniforme
di $\{0,\ldots,m-1\}$, e le sequenze per chiavi diverse sono indipendenti.
Si assume inoltre assenza di cancellazioni (i marcatori \texttt{DELETE}
degradano le prestazioni).

\begin{table}[ht]
	\centering
	\caption{Numero medio di ispezioni con indirizzamento aperto}
	\begin{tabular}{lc}
		\toprule
		\textbf{Operazione}    & \textbf{Caso medio}                      \\
		\midrule
		Ricerca con successo   & $\frac{1}{\alpha}\ln\frac{1}{1-\alpha}$  \\[6pt]
		Ricerca con insuccesso & $\frac{1}{1-\alpha}$                     \\
		\bottomrule
	\end{tabular}
\end{table}

\begin{esempio}[Ispezioni in funzione di $\alpha$]
	Con $\alpha = 1/2$ (tabella riempita al 50\%):
	ricerca con insuccesso $\to 1/(1-1/2) = 2$ ispezioni;
	ricerca con successo $\to 2\ln 2 \approx 1.387$ ispezioni.
	
	Con $\alpha = 0.75$ (tabella riempita al 75\%):
	ricerca con insuccesso $\to 4$ ispezioni;
	ricerca con successo $\to \frac{4}{3}\ln 4 \approx 1.848$ ispezioni.
\end{esempio}

% ............................................................
\subsubsection{Ispezione lineare}
% ............................................................

\begin{definizione}[Ispezione lineare]
	\[
	h(k, i) = \bigl(h_0(k) + i\bigr) \bmod m, \qquad i = 0, 1, \ldots, m-1
	\]
	In caso di collisione si scansionano le celle una ad una a partire da
	quella adiacente alla posizione originale.
\end{definizione}

\begin{hardnote}
	L'ispezione lineare soffre del \textbf{problema del clustering primario}:
	le sequenze di celle adiacenti occupate (\emph{cluster}) tendono a crescere
	rapidamente, vanificando la distribuzione uniforme. Il metodo genera al più
	$m$ permutazioni distinte invece di $m!$, quindi non soddisfa l'ipotesi di
	hashing con permutazioni indipendenti. È comunque utile su architetture con
	memoria gerarchica (cache/RAM/disco) grazie alla località spaziale degli
	accessi.
\end{hardnote}

% ............................................................
\subsubsection{Doppio hashing}
% ............................................................

\begin{definizione}[Doppio hashing]
	\[
	h(k, i) = \bigl(h_0(k) + i \cdot p(k)\bigr) \bmod m,
	\qquad i = 0, 1, \ldots, m-1
	\]
	dove $p(k)$ è una seconda funzione hash (\emph{funzione probe}). Affinché
	la sequenza copra tutte le $m$ celle è sufficiente che
	$\gcd\bigl(p(k), m\bigr) = 1$.
\end{definizione}

\begin{nota}
	Due costruzioni che garantiscono la copertura completa:
	\begin{itemize}
		\item $m$ primo e $0 < p(k) < m$: ad esempio $m = 701$,
		$p(k) = 1 + (k \bmod 500)$.
		\item $m = 2^c$ e $p(k)$ dispari: ad esempio $m = 512$,
		$p(k) = 1 + 2\,(k \bmod (m/2))$.
	\end{itemize}
	Il doppio hashing non soffre del clustering primario ed è in generale
	migliore dell'ispezione lineare e quadratica. Pur non soddisfacendo
	l'ipotesi di hashing con permutazioni indipendenti, le permutazioni
	prodotte ne hanno molte caratteristiche.
\end{nota}